% F08：本书原创矢量示意；依据所在节的定义、证明或明确模型。
\begin{figure}[htbp]
\centering
\begin{tikzpicture}[x=1cm,y=1cm,>=stealth,every node/.style={font=\small},line width=.65pt]

\foreach \x/\y/\r in {1/1/.6,1.8/1.2/.4,2.9/1/.55,3.5/1.5/.35,1.7/1.35/.3}{\draw[gray](\x,\y)circle(\r);}
\draw[AnalysisBlue,very thick](1,1)circle(.6);\draw[AnalysisBlue,very thick](2.9,1)circle(.55);
\node at(2.1,3.25){不交的选中球};
\draw[->](4.4,1.2)--(5.7,1.2);
\draw[AnalysisBlue,thick](8.7,1.25)circle(1.65);
\draw[AnalysisBlue,thick](8.7,1.25)circle(.33);
\draw[orange!80!black,thick](9.2,1.4)circle(.25);
\draw[dashed](8.7,1.25)--(9.2,1.4);
\node at(8.7,3.25){\(B'\subset 5B\)};
\node at(8.7,-1){舍弃球与所选球相交，且半径受控};

\end{tikzpicture}
\caption[5r 选球与膨胀覆盖]{5r 选球与膨胀覆盖。左图为有限选球示意；右图把一个选中球放大显示，外圆半径是内圆的五倍。覆盖结论仍依赖证明中的选取顺序与半径比较。}
\label{b1:fig:visual-f08}
\end{figure}
